\chapter{Lambda Calculus, Category Theory and Type Theory as it Pertains to Compilers}
\label{ch:app-d-foundations}

This appendix collects the small core of lambda calculus, category theory, and type theory that compiler work actually uses, in one place and in the notation of this book.
It introduces no new language feature and proves no new theorem about the book's passes.
A reader who can read a beta step, a typing rule, and a commuting square can follow every reference the chapters make to these ideas; a reader who wants the full theories should treat this appendix as a map, not a substitute.
% TODO(cite): need primary references for untyped lambda calculus, simply typed lambda calculus, and the categorical reading of types.

\section{Functions with nothing else}
\label{sec:app-d-lambda}

The \term{lambda calculus} is a language of functions with no numbers, no loops, and no machine.
Its terms have three shapes:
\begin{equation}
  e ::= x \mid \operatorname{lam}(x, e) \mid \operatorname{app}(e, e).
  \label{eq:app-d-syntax}
\end{equation}
The binder $\operatorname{lam}(x, e)$ names its parameter $x$ inside body $e$; free variables are those not bound by any enclosing binder.
Computation is \term{beta reduction}: applying an abstraction substitutes the argument for the bound parameter in the body.
Substitution must avoid capture: a free variable of the argument must not fall under a binder of the same spelling inside the body, so bound names are renamed first.
Algorithm~\ref{alg:app-d-subst} states the capture-avoiding step consumed below.
\begin{algorithm}
\caption{Capture-avoiding substitution of $N$ for $x$ in $M$.}
\label{alg:app-d-subst}
\begin{pseudocode}
\Require terms $M$ and $N$, variable $x$; bound names in $M$ distinct from free names of $N$ after renaming
\Ensure term $M[x := N]$ with no captured free variable of $N$
\Function{Subst}{$M, x, N$}
\If{$M$ is a variable $y$}
\If{$y = x$} \Return $N$ \Else \Return $y$ \EndIf
\ElsIf{$M$ is $\operatorname{app}(P, Q)$}
\State \Return $\operatorname{app}(\Call{Subst}{P, x, N}, \Call{Subst}{Q, x, N})$
\ElsIf{$M$ is $\operatorname{lam}(y, P)$ with $y = x$} \Return $M$
\Else \Comment{rename $y$ when $y$ occurs free in $N$}
\State pick fresh $z$ and \Return $\operatorname{lam}(z, \Call{Subst}{P[y := z], x, N})$
\EndIf
\EndFunction
\end{pseudocode}
\end{algorithm}
The freshness side condition is the whole content: without it, substitution confuses distinct variables that happen to share a spelling.

Numbers can be encoded as iteration counts: $\operatorname{zero}$ takes a step function and a base value and returns the base, while $\operatorname{succ}$ applies one more step.
Addition then composes iteration counts, and the worked step below evaluates with the book's own arithmetic rather than asserting it.
Consider $\operatorname{app}(\operatorname{lam}(x, x), \operatorname{lam}(y, y))$.
The argument is already a value, the binder $x$ occurs once, and no renaming is needed, so one beta step yields $\operatorname{lam}(y, y)$.
Figure~\ref{fig:app-d-beta} draws this step as an edge between two term nodes with the redex marked.
\begin{figure}
\centering
\begin{tikzpicture}[node distance=2.4cm]
\node[basic block] (redex) {app(lam(x, x), lam(y, y))};
\node[basic block, right of=redex, xshift=2.2cm] (value) {lam(y, y)};
\draw[cfg edge] (redex) -- node[above, font=\sffamily\scriptsize] {beta} (value);
\end{tikzpicture}
\caption{One beta step contracts the marked application and returns the argument unchanged; no capture is possible because the bound names differ.}
\label{fig:app-d-beta}
\end{figure}
A compiler rarely executes such encodings directly.
It uses them as a laboratory: evaluation order, scope extrusion, and capture discipline all appear here before they appear in closure conversion.

\section{Arrows keep application honest}
\label{sec:app-d-arrows}

The \term{simply typed lambda calculus} adds the arrow discipline of Chapter~\ref{ch:type-theory-compilers} to the calculus above.
Every abstraction carries its domain annotation, every application is checked against it, and Table~\ref{tab:ch39-rules} is the complete checker for the pure fragment.
A closed, well-typed term of ground kind evaluates to a value of that kind without confusing functions with base values; this is the kind-safety claim the chapter states for Mini-Fun, restricted here to the loop-free fragment.
Adding a fixed-point combinator or general recursion restores nontermination and must be recorded as an explicit extension with its own guard, not smuggled in as an example.
Type checking for this fragment is decidable and syntax-directed: Algorithm~\ref{alg:ch39-check} visits each node once under its context.
Polymorphism, existentials, and dependency extend the same pattern with instantiation, hiding, and term-indexed kinds, as surveyed in Section~\ref{sec:ch39-richer}; each needs its own solver and is therefore named but not developed here.

\section{Categories as an organizing analogy}
\label{sec:app-d-categories}

A \term{category} has objects, arrows between objects, an identity arrow on each object, and associative composition with identity laws.
The compiler-relevant reading is deliberately thin: take types as objects and well-typed programs as arrows, so composition is sequencing and identity is the do-nothing program.
A \term{product} then packages two values with projections that recover each side, which is the categorical shape of a pair; dually, a \term{coproduct} offers two injections with a case analysis that consumes either side, which is the shape of a tagged variant.
A \term{functor} maps one such structure into another while preserving identity and composition, and a \term{monad} packages an effect discipline (such as partiality, state, or exceptions) with a return and a sequencing operation obeying three laws.
This book uses only the first half of that sentence: products explain pairs, coproducts explain branching, and monads name an interface boundary where an optimization must not reorder effects.
No pass in this book proves an optimization by a universal property, and no chapter assumes that every well-typed program denotes a categorical arrow in a single category.
Readers who want that development should consult a dedicated text; here the vocabulary exists so that a phrase like \enquote{the product rule} or \enquote{the monadic sequencing point} points at one precise shape.

Table~\ref{tab:app-d-correspondence} records the three-way correspondence at the level the book actually uses.
\begin{table}
\centering
\caption{The working correspondence used by the book: each row names one shape in three vocabularies, without claiming the full theories coincide.}
\label{tab:app-d-correspondence}
\begin{tabular}{lll}
\toprule
Shape & Types reading & Categorical reading \\
\midrule
Function & $T_1 \to T_2$ with Lam, App & exponential with apply, curry \\
Pair & constructor, two projections & product with pairing, projections \\
Choice & one of two kinds, case analysis & coproduct with injections, cotupling \\
Effect boundary & sequencing point for effects & monadic bind respecting order \\
\bottomrule
\end{tabular}
\end{table}
The table is an analogy with exact edges, not an identification of the theories.
A new type feature earns a row only with its introduction rule, its elimination rule, and the equation that relates them; a new categorical claim earns its place only with the universal property written out.
% TODO(cite): need primary references for the Curry-Howard correspondence and the categorical semantics of simple types.

\section{What compilers take from these theories}
\label{sec:app-d-compiler-use}

Three borrowings recur across the book.
First, intermediate forms such as A-normal form and continuation-passing style make evaluation order and control joins explicit; they are lambda-calculus ideas restricted to a compiler's needs, developed in Chapter~\ref{ch:intermediate-representations} without importing the whole theory.
Second, closure conversion turns a nested function into a code pointer plus an environment holding exactly the free variables that substitution would have carried; the correctness question is which variables are free and where capture would occur, which is Section~\ref{sec:app-d-lambda} applied to Chapter~\ref{ch:inlining-specialization}.
Third, typed intermediate languages carry the judgments of Chapter~\ref{ch:type-theory-compilers} through lowering so that each pass checks its own output; the survey in Section~\ref{sec:ch06-ir-families} records where that discipline pays and where it costs.
In each case the theory supplies the shape of the invariant and the compiler supplies the data structures, the complexity bound, and the tests.
Confusion runs one way: importing vocabulary without the invariant produces passes that name a property they do not maintain.
The chapters therefore cite a rule or a map only where the pass enforces it, and mark every stronger claim with a TODO rather than asserting it.
